% =========================================================
\section{Молекулярно-орбітальне (МО) представлення}
% =========================================================

Хоча рівняння Хартрі–Фока розв’язуються в базисі атомних орбіталей $\{\chi_\mu\}$ (АО-представлення), після розв’язання ми завжди виконуємо перехід до \textbf{молекулярно-орбітального (МО) представлення}. У ньому одноелектронні хвильові функції — це вже не атомні орбіталі, а їх лінійні комбінації — \textbf{молекулярні орбіталі} $\{\varphi_p\}$.

Перехід від АО до МО здійснюється за допомогою матриці коефіцієнтів $\mathbf{C}$ (розмірність $M\times M$, де $M$ — число базисних функцій):

\begin{equation}\label{eq:MO_expansion}
    \varphi_p(\mathbf{r}) = \sum_{\mu=1}^{M} C_{\mu p}\,\chi_\mu(\mathbf{r}),
    \qquad p = 1,2,\dots,M
\end{equation}
або в компактній матричній формі:
\begin{equation}
    \boldsymbol{\varphi} = \boldsymbol{\chi}\,\mathbf{C}.
\end{equation}

Стовпці матриці $\mathbf{C}$ — це і є шукані молекулярні орбіталі, виражені в базисі атомних орбіталей.

Перехід до МО-представлення необхідний з кількох ключових причин:

\begin{enumerate}
    \item \textbf{Діагоналізація фокіана.}
          У канонічному МО-базисі матриця Фока стає діагональною:
          $\mathbf{F}^{\text{MO}} = \boldsymbol{\varepsilon} = \mathrm{diag}(\varepsilon_1,\varepsilon_2,\dots,\varepsilon_M)$.
    \item \textbf{Чіткий поділ на зайняті та віртуальні орбіталі.}
          Матриця густини в МО-базисі має блочно-діагональний вигляд: зайняті орбіталі — з населенням 2 (RHF) або 1 (UHF/ROHF), віртуальні — 0.
    \item \textbf{Зручність для пост-HF методів.}
          Усі методи кореляції (MP2, CI, CC, CASSCF), а також лінійний відгук (TDHF, TDDFT, CPHF) природно формулюються саме в МО-базисі. Двоелектронні інтеграли набувають компактної антисиметризованої форми $(ij\|ab)$.
\end{enumerate}

Таким чином, МО-представлення — це не просто технічна заміна базису, а основа сучасної інтерпретації результатів HF-розрахунків і фундамент для всіх наступних рівнів квантово-хімічної теорії.

\begin{commentbox}[Заміна базису: АО $\to$ МО]
Перехід від АО- до МО-представлення — це стандартна унітарна трансформація.
Якщо $\mathbf{C}$ --- матриця переходу (стовпці --- МО-коефіцієнти), то:

\begin{equation*}
  \mathbf{O}^{\text{MO}} = \mathbf{C}^\dagger \mathbf{O} \mathbf{C}.
\end{equation*}
\end{commentbox}

% ---------------------------------------------------------
\subsubsection{Матриця густини в МО-представленні}
% ---------------------------------------------------------

Матриця густини в атомно-орбітальному (АО) базисі для RHF має вигляд:
\begin{equation}
    \mathbf{D}^{\text{AO}} = 2 \sum_{i=1}^{n_{\text{occ}}} \mathbf{C}_i \mathbf{C}_i^\dagger
    = 2 \, \mathbf{C}_{\text{occ}} \mathbf{C}_{\text{occ}}^\dagger,
\end{equation}
де $\mathbf{C}_{\text{occ}}$ — підматриця коефіцієнтів зайнятих молекулярних орбіталей,
а множник 2 відповідає подвійному заповненню (спіни $\alpha$ і $\beta$).

Для переходу в МО-базис необхідно врахувати, що МО-коефіцієнти $\mathbf{C}$
ортонормовані \textbf{в метриці матриці перекриття} $\mathbf{S}$:
\begin{equation}
    \mathbf{C}^\dagger \mathbf{S} \mathbf{C} = \mathbf{I}.
\end{equation}

Трансформація матриці густини в МО-базис виконується через:
\begin{equation}\label{eq:D_MO_transform}
    \mathbf{D}^{\text{MO}} = \mathbf{C}^\dagger \mathbf{S} \, \mathbf{D}^{\text{AO}} \, \mathbf{S} \mathbf{C}.
\end{equation}

Для канонічних молекулярних орбіталей (власних векторів оператора Фока)
матриця густини у МО-представленні є \textbf{діагональною}:
\begin{equation}\label{eq:D_MO_diagonal}
    D_{pq}^{\mathrm{MO}} = n_p\,\delta_{pq},
\end{equation}
де $n_p$ --- заселеність орбіталі $p$:
\begin{equation}
    n_p =
    \begin{cases}
        2, & \text{RHF, зайнята орбіталь},\\[4pt]
        1, & \text{UHF/ROHF, зайнята $\alpha$ або $\beta$ орбіталь},\\[4pt]
        0, & \text{віртуальна орбіталь}.
    \end{cases}
\end{equation}

\textbf{Важливо:} пряме перетворення $\mathbf{C}^\dagger \mathbf{D}^{\text{AO}} \mathbf{C}$
\textit{без} матриці перекриття $\mathbf{S}$ є некоректним і не дає діагональної матриці.


% ---------------------------------------------------------
\subsubsection{Фокіан і одноелектронні оператори}
% ---------------------------------------------------------

Фокіан у МО-базисі:

\begin{equation}
    \mathbf{F}^{\text{MO}} = \mathbf{C}^\dagger \mathbf{F} \mathbf{C}
    = \boldsymbol{\varepsilon}
    = \mathrm{diag}(\varepsilon_1,\varepsilon_2,\dots,\varepsilon_M).
\end{equation}

Аналогічно перетворюються будь-які одноелектронні оператори (наприклад, дипольний момент, кінетична енергія тощо):

\begin{equation}
    \mathbf{h}^{\text{MO}} = \mathbf{C}^\dagger \mathbf{h} \mathbf{C}.
\end{equation}

% ---------------------------------------------------------
\subsubsection{Двоелектронні інтеграли в МО-базисі}
% ---------------------------------------------------------

Двоелектронні інтеграли в МО-представленні отримують шляхом чотирикратної трансформації:

\begin{equation}\label{eq:MO_eri}
    (pq|rs)
    = \sum_{\mu\nu\lambda\sigma}
      C_{\mu p} C_{\nu q} C_{\lambda r} C_{\sigma s}\,
      (\mu\nu|\lambda\sigma).
\end{equation}

У пост-HF методах часто використовується антисиметризована форма $(ij\|ab)$.

Наступний код демонструє розрахунки в МО-базисі:
\begin{minted}{python}
from pyscf import gto, scf
import numpy as np

mol = gto.M(atom='H 0 0 0; H 0 0 1.0', basis='sto-3g', verbose=0)
mf = scf.RHF(mol).run()

# --- матриця C_{μp} ---
C = mf.mo_coeff

# --- матриця густини ---
D = mf.make_rdm1()               # густина в АО

# --- Матриця Фока ---
F = mf.get_fock()

# --- матриця перекриття ---

S = mol.intor('int1e_ovlp')
print("\nМатриця густини в МО-базисі (правильна):")
D_MO = C.T @ S @ D @ S @ C
print(np.round(D_MO, 6))

# --- Матриця Фока в МО-основі (повинна бути діагональною!) ---
F_MO = C.T @ F @ C
print("\nМатриця Фока в МО-представлені:\n")
print(np.round(F_MO, 6))

# --- Орбітальні енергії з PySCF ---
eps = mf.mo_energy
print("\nОрбітальні енергії:\n")
print(np.round(eps, 6))
\end{minted}
% ---------------------------------------------------------
\subsubsection{Енергія Хартрі–Фока в МО-представленні}
% ---------------------------------------------------------

У МО-базисі вираз для повної енергії набуває особливо прозорого вигляду (для RHF):

\begin{align}\label{eq:E_HF_MO_explicit}
    E_{\mathrm{HF}}
    &= \sum_{i}^{\text{occ}} 2\,\varepsilon_i
       - \frac{1}{2} \sum_{i,j}^{\text{occ}} \bigl(2(ij|ij) - (ij|ji)\bigr)
       + V_{NN} \notag \\[6pt]
    &= \sum_{i}^{\text{occ}} 2\,\varepsilon_i
       - \sum_{i,j}^{\text{occ}} \bigl((ij|ij) - \tfrac{1}{2}(ij|ji)\bigr)
       + V_{NN}.
\end{align}

Альтернативна (часто зручніша) форма:

\begin{equation}
    E_{\mathrm{HF}}
    = \sum_{i}^{\text{occ}} 2\,\varepsilon_i
      - \frac{1}{2} \sum_{i,j}^{\text{occ}} (2J_{ij} - K_{ij})
      + V_{NN}.
\end{equation}

% ---------------------------------------------------------
\subsubsection{МО-формулювання рівнянь Хартрі–Фока}
% ---------------------------------------------------------

Вихідна матрична форма в АО-базисі:

\begin{equation}
    \mathbf{F} \mathbf{C} = \mathbf{S} \mathbf{C} \boldsymbol{\varepsilon}
\end{equation}

після переходу до МО-представлення зводиться до тривіального вигляду:

\begin{equation}
    \mathbf{F}^{\text{MO}} = \boldsymbol{\varepsilon}.
\end{equation}

Саме тому канонічні молекулярні орбіталі — це власні функції ефективного одноелектронного гамільтоніана (фокіана) в середньому полі інших електронів.